% Chapter 11: Sample Complexity, Compression, and Occam's Razor
% sample_complexity = Profile B-light, unlabeled_compression = Profile D (Negative),
% occam_algorithm = Profile B-light, description_length/kolmogorov/kl = Profile A,
% covering_number = Profile A-plus.
%
% Texture: question-driven. The open conjecture (compression ≈ VC dimension)
% is the gravitational center. Everything else orbits it: sample complexity
% provides the motivation, Occam's razor provides the algorithmic bridge,
% information-theoretic notions provide the vocabulary. The chapter should
% feel like a detective story where the case remains unsolved.

\chapter{Sample Complexity, Compression, and Occam's Razor}
\label{ch:compression}

A learner that processes a thousand training examples may ultimately lean on
three.  The Fundamental Theorem of \Cref{ch:pac} tells us \emph{which} classes
are PAC learnable, but it says nothing about what the learner actually needs to
retain.  The sharper question is: \emph{how little of the data must survive
for learning to succeed?}

The tight sample complexity bounds of \Cref{sec:sample-bounds} establish that
the VC dimension $d$ determines how many examples are needed, up to constant
factors.  Compression schemes (\Cref{sec:compression}) then reframe the same
question from the learner's side: a class is learnable if and only if any
training sample can be reduced to a bounded subsample that fully determines the
hypothesis.  The size of that subsample is what the chapter is really about.
The natural conjecture, that $O(d)$ points always suffice, was stated by
Littlestone and Warmuth in 1986.  It has not been resolved in the forty years
since.  Every known learnable class compresses to $O(d)$ points; no one can
prove that all of them do.

Two landmarks bracket the frontier.  Moran and Yehudayoff proved in 2016 that
compression to $2^{O(d)}$ points always exists: a fact unknown before their
work.  For $d = 20$, that bound is roughly $10^6$ points; the conjecture asks
for $20$.  That distance is the gap the conjecture has not closed in forty years.

Occam's razor (\Cref{sec:occam}) provides the bridge between compression and
generalization: any algorithm that finds a description of the data short enough
relative to the sample size generalizes.  The information-theoretic foundations
of \Cref{sec:info-foundations} (description length, Kolmogorov complexity,
covering numbers) give precise meaning to ``short enough.''

\medskip

The chapter is organized around a single question:

\begin{center}
\itshape
Does every concept class of VC dimension $d$ admit a compression scheme of
size $O(d)$?
\end{center}

Everything we prove either approaches this question or illuminates why it
resists resolution.


% ================================================================
% SECTION 1: TIGHT SAMPLE COMPLEXITY BOUNDS
% Profile B-light. State the bounds, give proof sketch for upper bound.
% ================================================================

\section{Tight Sample Complexity Bounds}
\label{sec:sample-bounds}

The sample complexity of PAC learning was bounded in \Cref{ch:pac} through the
VC characterization.  Here we state the tight bounds precisely.

\begin{thm}[Realizable PAC sample complexity {\cite{Hanneke2016,SimonZilles2015}}]
\label{thm:sample-realizable}
Let $\mathcal{H}$ be a hypothesis class with $\VC(\mathcal{H}) = d < \infty$.
In the realizable setting, the optimal sample complexity of PAC learning
$\mathcal{H}$ satisfies
\[
  m(\varepsilon, \delta) = \Theta\!\left(\frac{d}{\varepsilon}
    + \frac{\log(1/\delta)}{\varepsilon}\right).
\]
The upper bound is achieved by any consistent learner (including $\ERM$); the
lower bound holds for every learner.
\end{thm}

\begin{remark}[History of the tight bound]
\label{rmk:sample-history}
The classical upper bound of Blumer et al.~\cite{BlumerEhrenfeuchtHausslerWarmuth1989}
gave $O\!\bigl(\frac{d}{\varepsilon}\log\frac{1}{\varepsilon} +
\frac{1}{\varepsilon}\log\frac{1}{\delta}\bigr)$, with a logarithmic overhead
in the $d/\varepsilon$ term.  Removing this overhead (proving that the correct
dependence on $d$ is linear, not $d \log(1/\varepsilon)$) required substantial
effort.  The matching lower bound was established by Ehrenfeucht et
al.~\cite{EhrenfeuchtHausslerKearnsValiant1989}.  The tight upper bound without
the logarithmic factor was completed by Simon and Zilles~\cite{SimonZilles2015}
and, in a sharper form addressing the constant, by
Hanneke~\cite{Hanneke2016}.  Hanneke's result shows that \emph{any} consistent
learner achieves the optimal bound: no algorithmic cleverness beyond
consistency is needed.
\end{remark}

\begin{proof}[Proof sketch of the upper bound]
The argument proceeds in two stages.  First, one establishes the \emph{double
sampling} bound: the probability that $\ERM$ returns a hypothesis with error
$> \varepsilon$ is at most
\[
  \Pr[\exists h \in \mathcal{H} : \emprisk_S(h) = 0 \text{ and }
    \risk_D(h) > \varepsilon]
  \;\leq\; 2\,\growth_\mathcal{H}(2m) \cdot 2^{-\varepsilon m / 2},
\]
where $\growth_\mathcal{H}(n) \leq (en/d)^d$ by the Sauer--Shelah lemma
(\Cref{lem:sauer-shelah}).  Setting the right-hand side $\leq \delta$ and
solving for $m$ gives $m = O\!\bigl(\frac{d}{\varepsilon}
\log\frac{1}{\varepsilon} + \frac{1}{\varepsilon}\log\frac{1}{\delta}\bigr)$.

The removal of the $\log(1/\varepsilon)$ factor uses a more delicate argument.
One decomposes the hypothesis class into subsets of bounded ``disagreement mass''
and applies uniform convergence separately to each piece.  The key insight
(Hanneke~\cite{Hanneke2016}) is that the number of effectively distinct
hypotheses at accuracy scale $\varepsilon$ is controlled by the \emph{star
number} of $\mathcal{H}$, which is $O(d)$ rather than $O(d \log(1/\varepsilon))$.
\end{proof}

In the agnostic setting, the sample complexity worsens:

\begin{thm}[Agnostic PAC sample complexity]
\label{thm:sample-agnostic}
In the agnostic setting, the optimal sample complexity satisfies
\[
  m(\varepsilon, \delta) = \Theta\!\left(\frac{d}{\varepsilon^2}
    + \frac{\log(1/\delta)}{\varepsilon^2}\right).
\]
\end{thm}

The gap between $d/\varepsilon$ (realizable) and $d/\varepsilon^2$ (agnostic)
is the ``$\varepsilon^2$ price'' discussed in \Cref{sec:agnostic} of
\Cref{ch:pac}.  It is not an artifact of loose analysis: the lower bound
matches.


% ================================================================
% SECTION 2: COMPRESSION SCHEMES
% Profile F (Open Frontier). CENTERPIECE of the chapter.
% ================================================================

\section{Compression Schemes}
\label{sec:compression}

Sample complexity tells us how much data learning requires.  Compression asks
what happens once the data has arrived: \emph{how little of it does the learner
actually need to keep?}

\begin{defn}[Sample Compression Scheme {\cite{LittlestoneWarmuth1986,FloydWarmuth1995}}]
\label{def:compression-scheme}
A \emph{sample compression scheme} of size $k$ for a hypothesis class
$\mathcal{H}$ over domain $X$ consists of two maps:
\begin{enumerate}[leftmargin=*, itemsep=2pt]
\item A \emph{compression function} $\kappa$ that, given any sample $S$
  consistent with some $h \in \mathcal{H}$, selects a subsample
  $\kappa(S) \subseteq S$ with $|\kappa(S)| \leq k$.
\item A \emph{reconstruction function} $\rho$ that, given $\kappa(S)$, produces
  a hypothesis $\rho(\kappa(S))$ consistent with $S$.
\end{enumerate}
The reconstruction function $\rho$ sees \emph{only} the compressed subsample
$\kappa(S)$; it has no access to the remaining points.\formal{FLT_Proofs.Complexity.Structures.CompressionScheme}
\end{defn}

\begin{example}[Support vectors as compression]
\label{ex:svm-compression}
A support vector machine in $\R^d$ with a hard margin computes a
hyperplane determined by at most $d+1$ support vectors.  These support vectors
form a compression scheme of size $d+1$: the subsample $\kappa(S)$ consists of
the support vectors, and the reconstruction function $\rho$ returns the
maximum-margin hyperplane through them.
\end{example}

The fundamental connection between compression and learnability was established
by Littlestone and Warmuth:

\begin{thm}[Compression implies learnability {\cite{LittlestoneWarmuth1986,FloydWarmuth1995}}]
\label{thm:compression-pac}
If $\mathcal{H}$ admits a sample compression scheme of size $k$, then
$\mathcal{H}$ is PAC learnable with sample complexity
\[
  m(\varepsilon, \delta) = O\!\left(\frac{k}{\varepsilon}
    \log\frac{k}{\varepsilon} + \frac{1}{\varepsilon}\log\frac{1}{\delta}\right).
\]
In particular a finite compression scheme forces finite VC dimension, $\VC(\mathcal{H}) \leq 2(k+1)^2 - 1$, so compressibility implies learnability.\formal{FLT_Proofs.Complexity.Generalization.compression_imp_vcdim_finite}
\end{thm}

\begin{proof}[Proof sketch]
Fix a sample $S$ of size $m$ drawn i.i.d.\ from $D$.  The compression function
selects a subsample $\kappa(S)$ of size $\leq k$.  The number of possible
subsamples is $\binom{m}{k} \leq m^k$.  For each fixed subsample, the
reconstruction $\rho(\kappa(S))$ is a fixed hypothesis, and by a Hoeffding
bound, a single fixed hypothesis with empirical risk zero on the remaining
$m - k$ points has true risk $> \varepsilon$ with probability at most
$e^{-\varepsilon(m-k)}$.  Taking a union bound over all $m^k$ possible
subsamples:
\[
  \Pr[\risk_D(\rho(\kappa(S))) > \varepsilon] \leq m^k \cdot e^{-\varepsilon(m-k)}.
\]
Setting this $\leq \delta$ and solving for $m$ gives the stated bound.
\end{proof}

The converse direction (does learnability imply compression?) is where the
story becomes a forty-year open problem.

\subsection{The Compression Conjecture}

\begin{openprob}[Sample Compression Conjecture {\cite{LittlestoneWarmuth1986,FloydWarmuth1995}}]
\label{op:compression}
Does every concept class $\mathcal{H}$ with $\VC(\mathcal{H}) = d$ admit a
sample compression scheme of size $O(d)$?\formal{FLT_Proofs.Complexity.Structures.Open_NoInfoCompressionStrengthening}
\end{openprob}

This conjecture, stated by Littlestone and Warmuth in 1986, remains open.  It
asks for the converse of \Cref{thm:compression-pac}: if a class is learnable
(which, by the Fundamental Theorem, is equivalent to finite VC dimension $d$),
can the learner always compress its evidence to $O(d)$ points?

Four decades on, the statement is clean, the motivation immediate, and every
natural proof strategy fails.  Topological arguments reach $2^{O(d)}$ but not
$O(d)$.  Algebraic approaches founder on classes with no linear structure.  The
gap between what compression \emph{exists} and what the conjecture \emph{demands}
has not narrowed since the question was posed.

\begin{thm}[Exponential compression exists {\cite{MoranYehudayoff2016}}]
\label{thm:moran-yehudayoff}
Every concept class $\mathcal{H}$ with $\VC(\mathcal{H}) = d$ admits a sample
compression scheme of size $2^{O(d)}$.\formal{FLT_Proofs.Theorem.PAC.fundamental_vc_compression}
\end{thm}

This result of Moran and Yehudayoff~\cite{MoranYehudayoff2016}
established that compression schemes \emph{exist} for every learnable class, a
fact unknown before 2016.  The exponential bound $2^{O(d)}$ is vastly larger
than the conjectured $O(d)$.  The scheme it produces carries side information;
whether the side bits can be removed, and whether the size can be brought to
$O(d)$, are the open strengthenings (\cref{op:compression}).

\begin{proof}[Proof sketch]
The proof uses a topological argument.  By a result of
Radon, any set of $d+1$ points in a $d$-dimensional space admits a Radon
partition.  Moran and Yehudayoff generalize this to VC classes via a
``staircase'' construction: they build a sequence of hypothesis classes of
increasing VC dimension and show that each admits a compression scheme, using
the Sauer--Shelah lemma (\Cref{lem:sauer-shelah}) to control the growth of
the compression size.  The exponential blowup arises from iterating this
construction $d$ times.
\end{proof}

\begin{remark}[The gap]
\label{rmk:compression-gap}
The gap between the conjecture ($O(d)$) and the best known upper bound
($2^{O(d)}$) is enormous.  For $d = 20$, the conjecture predicts compression to
$\sim 20$ points; the Moran--Yehudayoff bound gives $\sim 10^6$.  No
polynomial-in-$d$ bound is known for general classes.
\end{remark}

\subsection{The One-Inclusion Graph}
\label{sec:one-inclusion}

The Moran--Yehudayoff proof rests on a combinatorial object that deserves
explicit treatment, both because it illuminates why maximum classes are
special and because it localizes the difficulty of the conjecture at a
precise structural point.

\begin{defn}[One-Inclusion Graph {\cite{HausslerLittlestonWarmuth1994}}]
\label{def:one-inclusion-graph}
Let $\mathcal{H}$ be a concept class over domain $X$, and let
$S = (x_1, \ldots, x_m)$ be a sequence of points from $X$.  The
\emph{one-inclusion graph} $G(\mathcal{H}, S)$ is the graph whose
vertex set is
$\{(h(x_1), \ldots, h(x_m)) : h \in \mathcal{H}\}$
(the distinct restrictions of $\mathcal{H}$ to $S$)
and whose edge set connects pairs of restrictions that differ on
exactly one coordinate.
\end{defn}

The one-inclusion graph mediates between the VC dimension and
compression.  The key property is that its density (specifically, the
minimum over all orientations of the maximum in-degree) controls
whether a compression scheme of size $d$ exists.

\begin{thm}[One-Inclusion Graph Orientation {\cite{RubinsteinKroese2007,KuzminWarmuth2007}}]
\label{thm:one-inclusion-orientation}
Let $\mathcal{H}$ be a \emph{maximum class} of VC dimension $d$
(i.e., $\growth_\mathcal{H}(m) = \sum_{i=0}^d \binom{m}{i}$ for all
$m$).  Then for any sample $S$ of size $m \geq d$, the one-inclusion
graph $G(\mathcal{H}, S)$ can be oriented so that every vertex has
in-degree at most $d$.
\end{thm}

This orientation directly yields a compression scheme of size $d$.
Given a sample $S$ consistent with concept $h^*$, the vertex
$v = h^*|_S$ in the oriented one-inclusion graph has at most $d$
incoming edges.  Each incoming edge corresponds to a sample point
$x_i$ where some neighboring concept $h'$ disagrees with $h^*$.
The compressed subsample $\kappa(S)$ consists of these $\leq d$
``witness points.''  The reconstruction function $\rho$ recovers
$h^*$ from the witness points using the acyclicity of the oriented
graph: the witness points uniquely determine $h^*$ among all
concepts in $\mathcal{H}$ consistent with $S$.

\begin{example}[One-inclusion graph for intervals]
\label{ex:oig-intervals}
Let $\mathcal{H}$ be the class of intervals $[a, b] \subseteq [0, 1]$
(labeling $x$ as $1$ iff $x \in [a,b]$), with $\VC(\mathcal{H}) = 2$.
This is a maximum class.  Given a sorted sample
$S = (x_1 < x_2 < \cdots < x_m)$, each concept restricts to a binary
string of the form $0^i 1^j 0^k$ (a contiguous block of 1s flanked by
0s).  The one-inclusion graph connects patterns that differ in a single
bit, which can only happen at the boundaries of the block: either the
leftmost $1$ becomes $0$ or the rightmost $1$ becomes $0$ (or the
reverse).  The graph can be oriented so that each vertex points toward
its two boundary points.  The compressed subsample consists of the two
boundary points of the interval: exactly the information needed to
reconstruct $[a, b]$ on the sample.

The mechanism is subtle: what matters is the two \emph{sample points}
at the boundary of the positive region.  The \emph{endpoints} of the
interval in $[0,1]$ play no role (and recovering them would require
real-valued parameters, which the scheme never stores).  The
one-inclusion graph identifies the two boundary sample points
combinatorially, without reference to the underlying geometry.
\end{example}

For general (non-maximum) classes, the one-inclusion graph does
\emph{not} admit an orientation with in-degree $\leq d$.  This is
the precise technical point where the conjecture resists.

\begin{remark}[The maximum-to-general gap]
\label{rmk:max-general-gap}
Moran and Yehudayoff's strategy is to reduce the general case to the
maximum case.  Any class $\mathcal{H}$ of VC dimension $d$ can be
``fattened'' to a maximum class $\mathcal{M}$ with
$\VC(\mathcal{M}) = d$ by adding concepts until the Sauer--Shelah
bound is achieved with equality.  The maximum class has compression of
size $d$, and the compression scheme can in principle be pulled back to
$\mathcal{H}$.  However, the fattening introduces auxiliary
concepts; the reconstruction function, given a compressed subsample,
must decide which of many $\mathcal{M}$-consistent concepts actually
belongs to $\mathcal{H}$.  The number of such choices is exponential
in $d$, and encoding this choice inflates the compression size to
$2^{O(d)}$.

Closing the gap to $O(d)$ requires either working directly with
non-maximum classes (for which no orientation theorem is known) or
proving that the reconstruction ambiguity can be resolved with
$O(d)$ additional bits rather than $2^{O(d)}$.
\end{remark}


The conjecture has been verified for specific class families:

\begin{enumerate}[leftmargin=*, itemsep=3pt]
\item \textbf{Maximum classes.}  A class $\mathcal{H}$ with
  $\growth_\mathcal{H}(n) = \sum_{i=0}^d \binom{n}{i}$ for all $n$ (achieving
  the Sauer--Shelah bound with equality) admits a compression scheme of size
  $d$~\cite{FloydWarmuth1995,RubinsteinKroese2007}.

\item \textbf{Intersection-closed classes.}  If $\mathcal{H}$ is closed under
  intersections, compression of size $d$ exists.

\item \textbf{Classes of VC dimension 1.}  Trivially: a single point suffices.

\item \textbf{Halfspaces in $\R^d$.}  The support vector compression of
  \Cref{ex:svm-compression} gives size $d + 1$.

\item \textbf{Balls, rectangles, and other geometric classes.}  Known
  case-by-case.
\end{enumerate}

The intersection-closed case is worth examining in detail, because its
compression mechanism is fundamentally different from the SVM case and
reveals what structure makes compression ``easy.''

\begin{example}[Compression for intersection-closed classes]
\label{ex:intersection-closed-compression}
Let $\mathcal{H}$ be a concept class closed under intersection: if
$h_1, h_2 \in \mathcal{H}$, then $h_1 \cap h_2 \in \mathcal{H}$
(where $(h_1 \cap h_2)(x) = h_1(x) \wedge h_2(x)$).

Given a sample $S$ consistent with target $h^* \in \mathcal{H}$,
define
$\hat{h} = \bigcap\{h \in \mathcal{H} : h \text{ consistent with } S\}$.
By intersection-closure, $\hat{h} \in \mathcal{H}$, and $\hat{h}$ is
the \emph{unique minimal} concept consistent with~$S$.  The
compression scheme exploits this uniqueness.

Let $P^+ = \{x_i \in S : h^*(x_i) = 1\}$ be the positive examples.
The minimal concept $\hat{h}$ is determined by the constraint that it
must label all points in $P^+$ as positive while being minimal in
$\mathcal{H}$.  A \emph{spanning set} for $P^+$ is a subset
$T \subseteq P^+$ such that
$\bigcap\{h \in \mathcal{H} : h(x) = 1\;\forall x \in T\}
 = \bigcap\{h \in \mathcal{H} : h(x) = 1\;\forall x \in P^+\}$%
that is, $T$ imposes the same constraints as the full positive set.
A minimal spanning set has size $\leq d = \VC(\mathcal{H})$, because
the VC dimension bounds the number of independent constraints.  The
compression function selects a minimal spanning set; the
reconstruction function computes
$\hat{h} = \bigcap\{h \in \mathcal{H} : h(x) = 1\;\forall (x, 1) \in
\kappa(S)\}$.

The mechanism is structurally different from SVMs.  Support vector
machines use \emph{geometric} structure (margin maximization) to
select anchor points.  Intersection-closed classes use \emph{lattice}
structure (existence of meets) to guarantee a unique minimum.  Both
achieve compression of size $d$, but via incomparable mathematical
facts.  This incomparability is a symptom of the conjecture's
difficulty: there is no single mechanism that explains compression
across all class families.
\end{example}


\subsection{Compression for Attention Classes}
\label{sec:attention-compression}

The compression question is not a relic of classical concept classes; it
reappears, still unsolved, for the hypothesis classes computed by modern
attention layers.  Fix an embedding dimension $d$ and a context of $n_K$
key--value pairs.  A single softmax-attention head maps a query $x \in \R^d$ to
the weighted readout
$\sum_{i=1}^{n_K} \mathrm{softmax}(\langle x, k_i\rangle)_i\, v_i$, and
thresholding this readout at zero yields a concept
$c_{K,V}\colon \R^d \to \{0,1\}$.  Letting the keys $K = (k_1, \ldots, k_{n_K})$
and values $V = (v_1, \ldots, v_{n_K})$ range over all configurations gives the
\emph{attention concept class} $\mathcal{A}_{d,n_K} = \{c_{K,V}\}$.

\begin{prop}[Attention classes are compressible]
\label{prop:attention-compression}
For every $d$ and $n_K$, the attention concept class $\mathcal{A}_{d,n_K}$ has
finite VC dimension and therefore admits a sample compression scheme of size
$2^{O(\VC(\mathcal{A}_{d,n_K}))}$.
\end{prop}

\begin{proof}
The readout map
$(K, V, x) \mapsto \sum_i \mathrm{softmax}(\langle x,k_i\rangle)_i\, v_i$ is
jointly continuous, hence Borel measurable, so the bad-event sets the
uniform-convergence argument quantifies over are measurable.  This is the
well-behaved measurable VC target
condition,\formal{TLT.Bridge.softAttention_wellBehaved} so the empirical process
indexed by $\mathcal{A}_{d,n_K}$ converges uniformly and the class has finite VC
dimension.  Finite VC dimension is the hypothesis of the existence
theorem,\formal{FLT_Proofs.Theorem.PAC.fundamental_vc_compression} which supplies
a compression scheme of size exponential in the VC dimension.
\end{proof}

The exponential bound is, as for general classes, far from tight, and the argmax
(hard-attention) router suggests why.  When the softmax is sharpened to a top-$1$
selection, the head's output on a query $x$ is determined by the single key whose
score is largest, and the routing rule that selects this key is again
well-behaved,\formal{TLT_Proofs.Routing.MeasurableScoreRouting.multiHeadArgmax_wellBehaved}
partitioning the query space into finitely many Borel cells.  A query's label
depends only on which cell it lands in, and each cell is pinned down by a bounded
number of support keys (the attention analogue of the support vectors of
\Cref{ex:svm-compression}).  This predicts a compression of size $O(n_K)$, linear
in the context length.  Turning the heuristic into a theorem runs into the
obstacle of \Cref{op:compression}: the support-key set is selected with knowledge
of the whole sample, and pulling it back to a labeled subsample of size $O(\VC)$
is the open problem.

\begin{remark}[The conjecture, modernized]
\label{rmk:attention-compression}
\Cref{prop:attention-compression} places attention classes inside the landscape
of \Cref{fig:compression-landscape} at the compression-size $2^{O(d)}$ node, with
the descent to size $O(d)$ open.  The Littlestone--Warmuth conjecture is thus
load-bearing for contemporary architectures.  A positive resolution would bound
how few training examples an attention head must retain to reproduce its
behavior; a negative one would exhibit a learnable attention class that is
incompressible.
\end{remark}


\subsection{Why the Conjecture Resists Resolution}
\label{sec:why-hard}

The compression conjecture has resisted attack for four decades.  Techniques
are not scarce: the difficulty is structural.  Every known approach either
works for specific, well-structured class families but cannot reach arbitrary
VC classes, or works for all VC classes but pays an exponential price in scheme
size.  The two regimes do not overlap, and no method yet bridges them.

\begin{obstbox}
\textbf{Three approaches and their failure modes.}

\medskip
\emph{Approach 1: Direct construction.}  For each verified class
family (halfspaces, intervals, intersection-closed, maximum classes),
one can construct a compression scheme tailored to the family's
structure.  The constructions exploit distinct properties: geometric
anchoring for SVMs, lattice minimality for intersection-closed classes,
one-inclusion graph orientation for maximum classes.  But these
properties are incomparable: no single structural assumption covers all
families, let alone arbitrary VC classes.  Every direct construction
exploits something that arbitrary classes lack.

\medskip
\emph{Approach 2: One-inclusion graph orientation.}  For maximum
classes, \Cref{thm:one-inclusion-orientation} provides an orientation
with in-degree $d$, directly yielding compression of size $d$.  For
non-maximum classes, no such orientation theorem is known.
Non-maximum one-inclusion graphs can have vertices with degree much
larger than $d$, and no known orientation strategy achieves in-degree
$O(d)$ in general.  The maximum class case is tight, but the
combinatorial extremality that makes it work (achieving the
Sauer--Shelah bound with equality) is exactly what arbitrary classes
lack.

\medskip
\emph{Approach 3: Topological embedding.}  Moran and Yehudayoff's
proof uses a topological argument related to Radon partitions in
convex geometry.  The topological machinery is powerful enough to
handle all VC classes but pays an exponential price: the Radon-type
argument iterates a doubling construction $d$ times, producing
compression of size $2^{O(d)}$.  Improving this requires either a
fundamentally different topological tool or a way to avoid the
iterative blowup; the iterative structure appears intrinsic to
the proof method itself, the doubling having to run $d$ times before a
Radon partition is available, so a tighter analysis of the same
construction would not remove it.

\medskip
The conjecture sits at the intersection of three mathematical
territories: combinatorics (VC dimension, growth functions, maximum
classes), topology (Radon partitions, convex position), and
information theory (description length, reconstruction
complexity); and each territory contributes techniques that
illuminate a piece of the puzzle while leaving the central question
open.  A resolution may require connecting these territories in a way
that current methods do not.
\end{obstbox}


\subsection{Labeled vs.\ Unlabeled Compression}
\label{sec:unlabeled-compression}

In \Cref{def:compression-scheme}, the compressed subsample $\kappa(S)$ retains
both the points and their labels.  Labels carry information, and a natural
question is whether the labels are doing real work: can the reconstruction
function recover from the point positions alone?

\begin{defn}[Unlabeled Compression Scheme]
\label{def:unlabeled-compression}
An \emph{unlabeled compression scheme} of size $k$ for $\mathcal{H}$ is a pair
$(\kappa, \rho)$ where $\kappa$ selects a set of $\leq k$ \emph{points}
(without labels) from the sample, and $\rho$ reconstructs a hypothesis
consistent with $S$ from these points alone.
\end{defn}

The compression conjecture has an unlabeled analogue, in which the stored
subsample carries no labels and the reconstruction recovers them from the point
positions alone.

\begin{thm}[Bounded unlabeled compression {\cite{MoranYehudayoff2016}}]
\label{thm:unlabeled-bounded}
Every concept class $\mathcal{H}$ with $\VC(\mathcal{H}) = d$ admits an
unlabeled compression scheme of size bounded by a function of $d$ alone,
independent of $|\mathcal{H}|$.
\end{thm}

This matches the labeled case qualitatively: labeled compression of size
$2^{O(d)}$ always exists (\Cref{thm:moran-yehudayoff}), and the unlabeled scheme
stays bounded by a function of $d$ as well.  Whether the unlabeled size can be
pushed to $O(d)$, matching the labeled conjecture, is the Kuzmin--Warmuth
unlabeled compression conjecture~\cite{KuzminWarmuth2007}, open in general and
proved for maximum classes.

\begin{remark}
The unlabeled conjecture sharpens the labeled one (\Cref{op:compression}):
recovering labels from the point positions alone is the extra demand, and
whether it costs more than a constant factor in scheme size is unknown.
\end{remark}


% ================================================================
% SECTION 3: OCCAM'S RAZOR
% Profile B-light. Occam → PAC and partial converse.
% ================================================================

\section{Occam's Razor}
\label{sec:occam}

Compression schemes embody one form of Occam's razor: prefer the hypothesis
determined by the fewest data points.  A description-length version of the same
principle, formalized by Blumer, Ehrenfeucht, Haussler, and Warmuth, shifts the
criterion from subsample size to the length of the binary encoding: the output
hypothesis should be shorter than the sample itself.  Both forms connect
simplicity to generalization, and the formal connection runs precisely.

\begin{defn}[Occam Algorithm {\cite{BlumerEhrenfeuchtHausslerWarmuth1987}}]
\label{def:occam}
Let $\mathcal{H}$ be a hypothesis class and let $\mathcal{R}$ be a
\emph{representation class} (a set of hypotheses with associated description
lengths).  An algorithm $A$ is an \emph{Occam algorithm} for $\mathcal{H}$ with
respect to $\mathcal{R}$ if, given a sample $S$ of size $m$ consistent with some
$h \in \mathcal{H}$ of description length $n$, the algorithm outputs a
hypothesis $h' \in \mathcal{R}$ that:
\begin{enumerate}[leftmargin=*, nosep]
\item is consistent with $S$, and
\item has description length $|h'| \leq n^\alpha m^\beta$ for some constants
  $\alpha \geq 1$ and $0 \leq \beta < 1$.
\end{enumerate}
The sub-linear dependence on $m$ (exponent $\beta < 1$) is the key requirement:
the output hypothesis must be ``simpler'' than the data.
\end{defn}

\begin{thm}[Occam's Razor Theorem {\cite{BlumerEhrenfeuchtHausslerWarmuth1987}}]
\label{thm:occam}
If $A$ is an Occam algorithm for $\mathcal{H}$ with parameters $\alpha, \beta$,
then $A$ PAC-learns $\mathcal{H}$ with sample complexity
\[
  m(\varepsilon, \delta) = O\!\left(\frac{1}{\varepsilon}\left(
    n^\alpha \log\frac{1}{\varepsilon}
    + \log\frac{1}{\delta}\right)\right)^{1/(1-\beta)}.
\]
\end{thm}

\begin{proof}[Proof sketch]
The argument is a counting and union bound.  The number of hypotheses in
$\mathcal{R}$ with description length $\leq L$ is at most $2^L$.  An Occam
algorithm outputs a hypothesis of length $\leq n^\alpha m^\beta$.  Fixing a
sample of size $m$, the number of possible outputs is at most
$2^{n^\alpha m^\beta}$.  For each fixed hypothesis with empirical risk zero, the
probability that its true risk exceeds $\varepsilon$ is at most
$(1-\varepsilon)^m \leq e^{-\varepsilon m}$.  A union bound gives
\[
  \Pr[\risk_D(A(S)) > \varepsilon] \leq 2^{n^\alpha m^\beta} \cdot e^{-\varepsilon m}.
\]
Setting the right-hand side $\leq \delta$ and solving for $m$ (using $\beta < 1$
to ensure that the exponential decay in $m$ dominates the polynomial growth of
$2^{m^\beta}$) gives the stated bound.
\end{proof}

\begin{remark}[Occam and compression]
\label{rmk:occam-compression}
Every compression scheme of size $k$ induces an Occam algorithm: the compressed
subsample has description length $O(k \log m)$, which is sub-linear in $m$.
Conversely, an Occam algorithm is not necessarily a compression scheme, because
the output hypothesis need not be determined by a subsample of the input.
Occam's razor is thus a relaxation of compression.
\end{remark}

The converse of the Occam theorem (does every PAC-learnable class admit an
Occam algorithm?) was partially answered by Board and Pitt:

\begin{thm}[Partial converse {\cite{BoardPitt1992}}]
\label{thm:board-pitt}
If the class of all polynomial-size circuits is not PAC learnable by
polynomial-size circuits (a widely believed complexity-theoretic assumption),
then there exist PAC-learnable classes that do not admit polynomial-time Occam
algorithms.
\end{thm}

Thus the Occam property is \emph{strictly stronger} than PAC learnability
(under standard complexity assumptions): there exist learnable classes where
no efficient algorithm can find short descriptions of the data.  The gap is
purely computational: a short consistent description still exists, but the
work of producing it outruns any polynomial-time algorithm.


% ================================================================
% SECTION 4: INFORMATION-THEORETIC FOUNDATIONS
% Profile A and A-plus. Brief definitions.
% ================================================================

\section{Information-Theoretic Foundations}
\label{sec:info-foundations}

``Short description'' and ``efficiently coverable function class'' are the
informal language Occam's razor and compression schemes both rely on.  This
section gives each notion a precise, computable (or at least formal) definition.

\subsection{Description Length}

\begin{defn}[Description Length]
\label{def:description-length}
The \emph{description length} of a hypothesis $h$ in a representation class
$\mathcal{R}$ is the length $|h|$ of the shortest binary string encoding $h$ in
$\mathcal{R}$.  A \emph{prefix-free} representation ensures that
$\sum_{h \in \mathcal{R}} 2^{-|h|} \leq 1$ (the Kraft inequality).
\end{defn}

Description length connects to PAC learning through the observation that a
class of hypotheses with bounded description lengths has bounded effective size:
$|\{h \in \mathcal{R} : |h| \leq L\}| \leq 2^L$.  This makes union bounds
over the class tractable.

\subsection{Kolmogorov Complexity}

\begin{defn}[Kolmogorov Complexity]
\label{def:kolmogorov}
The \emph{Kolmogorov complexity} $K(x)$ of a string $x$ is the length of the
shortest program (on a fixed universal Turing machine) that outputs $x$ and
halts.  The \emph{conditional Kolmogorov complexity} $K(x \mid y)$ is the length
of the shortest program that outputs $x$ given $y$ as input.
\end{defn}

Kolmogorov complexity is uncomputable but provides the theoretical gold standard
for description length: $K(x)$ is the ultimate compression of $x$.

\subsection{KL Complexity}

\begin{defn}[KL Complexity]
\label{def:kl-complexity}
The \emph{Kullback--Leibler divergence} (KL divergence) between distributions
$P$ and $Q$ on a measurable space is
\[
  \KL(P \| Q) = \E_P\!\left[\log\frac{dP}{dQ}\right],
\]
when $P \ll Q$, and $+\infty$ otherwise.  In the PAC-Bayes framework
(\Cref{ch:generalization}), $\KL(\rho \| \pi)$ measures the ``complexity'' of a
posterior $\rho$ relative to a prior $\pi$: it is the number of additional
nats needed to encode a hypothesis drawn from $\rho$ using a code designed for
$\pi$.
\end{defn}

\subsection{Covering Numbers}

\begin{defn}[Covering Number]
\label{def:covering-number}
Let $(\mathcal{F}, d)$ be a (pseudo)metric space.  The \emph{$\varepsilon$-covering
number} $\mathcal{N}(\varepsilon, \mathcal{F}, d)$ is the minimum number of
balls of radius $\varepsilon$ (under $d$) needed to cover $\mathcal{F}$.
\end{defn}

Covering numbers measure the ``effective size'' of a function class at
resolution $\varepsilon$.  They connect to sample complexity through the
classical chaining argument:

\begin{thm}[Dudley's entropy integral {\cite{Dudley1967}}]
\label{thm:dudley}
Let $\mathcal{F}$ be a class of functions $f \colon X \to [0,1]$, and let
$d_n(f,g) = \bigl(\frac{1}{n}\sum_{i=1}^n (f(x_i) - g(x_i))^2\bigr)^{1/2}$
be the empirical $L^2$ metric on a sample $(x_1, \ldots, x_n)$.  Then
\[
  \E\!\left[\sup_{f \in \mathcal{F}} \left|\frac{1}{n}\sum_{i=1}^n f(x_i)
    - \E[f]\right|\right]
  \leq \frac{C}{\sqrt{n}} \int_0^1
    \sqrt{\log \mathcal{N}(\varepsilon, \mathcal{F}, d_n)}\; d\varepsilon,
\]
where $C$ is a universal constant.
\end{thm}

The integral $\int_0^1 \sqrt{\log \mathcal{N}(\varepsilon, \mathcal{F}, d_n)}\;
d\varepsilon$ is the \emph{Dudley entropy integral}.  It bounds the expected
supremum of the empirical process indexed by $\mathcal{F}$, the quantity that
controls uniform convergence and hence generalization.  When
$\mathcal{N}(\varepsilon, \mathcal{F}, d_n) \leq (C/\varepsilon)^d$ (as for
VC classes), the integral evaluates to $O(\sqrt{d})$, recovering the standard
sample complexity bounds.


% ================================================================
% SECTION 5: CHAPTER SUMMARY AND OPEN QUESTIONS
% ================================================================

\section{The State of the Art}
\label{sec:compression-summary}

The chapter's results cluster around a single qualitative fact and a single open
question.  The qualitative fact is settled and kernel-verified: PAC learnability,
finite VC dimension, and the existence of a finite compression scheme with side
information are equivalent.  The open question is how short the scheme can be
made.  \Cref{fig:compression-landscape} shows where each result sits and which
descent steps remain unproven.

\begin{figure}[ht]
\centering
\begin{tikzpicture}[
    box/.style={draw, thick, rounded corners, minimum width=2.5cm,
      minimum height=0.85cm, font=\small, align=center},
    verified/.style={box, fill=defblue!12},
    literature/.style={box, fill=defblue!4, draw=gray!55},
    openq/.style={box, dashed, fill=edgeorange!10, text=edgeorange!85!black},
    refuted/.style={box, fill=thmred!10, draw=thmred!70},
    veq/.style={{Stealth[length=2.5mm]}-{Stealth[length=2.5mm]}, thick, defblue!75},
    vedge/.style={-{Stealth[length=2.5mm]}, thick, defblue!75},
    litedge/.style={-{Stealth[length=2.5mm]}, semithick, gray!60},
    openedge/.style={-{Stealth[length=2.5mm]}, thick, dashed, edgeorange},
    rededge/.style={-{Stealth[length=2.5mm]}, thick, dashed, thmred},
    lbl/.style={font=\scriptsize, midway, fill=white, inner sep=1pt},
]
% --- verified equivalence core: THAT a short summary exists (kernel sc=0) ---
\node[verified] (pac) at (0,0) {PAC learnable};
\node[verified] (vc) at (3.5,0) {$\VC<\infty$};
\node[verified] (comp) at (7.0,0) {Finite\\compression\\{\scriptsize with side info}};
\draw[veq] (pac) -- (vc);
\draw[veq] (vc) -- (comp);

% --- Occam: a literature alternative route to the same qualitative conclusion ---
\node[literature] (occam) at (3.5,-2.0) {Occam algorithm};
\draw[litedge] (comp) -- node[lbl, above right] {implies} (occam);
\draw[litedge] (occam) -- node[lbl, below left] {\cref{thm:occam}} (pac);

% --- the size ladder: HOW short? descent = a strictly stronger claim ---
\node[literature] (sexp) at (7.0,-2.0) {Size $2^{O(d)}$};
\node[openq] (spoly) at (7.0,-3.4) {Size $\mathrm{poly}(d)$?};
\node[openq] (slin) at (7.0,-4.8) {Size $O(d)$?};
\draw[litedge] (comp) -- node[lbl, right] {\cite{MoranYehudayoff2016}} (sexp);
\draw[openedge] (sexp) -- node[lbl, right] {\cref{op:poly-compression}} (spoly);
\draw[openedge] (spoly) -- node[lbl, right] {\cref{op:compression}} (slin);

% --- side-info axis: an open strengthening of the verified core ---
\node[openq] (noinfo) at (10.1,-2.0) {Drop side info?};
\draw[openedge] (comp) -- node[lbl, pos=0.72] {\cref{op:side-info}} (noinfo);

% --- the unlabeled variant: bounded size known, O(d) open ---
\node[openq] (unlab) at (7.0,-6.2) {Unlabeled $O(d)$?};
\draw[openedge] (slin) -- node[lbl, right] {\cite{KuzminWarmuth2007}} (unlab);

% --- line-style legend: the evidence grade ---
\begin{scope}[shift={(0,-3.7)}]
  \draw[vedge] (0,0) -- (0.6,0);
  \node[font=\scriptsize, anchor=west] at (0.72,0) {kernel-verified ($sc{=}0$)};
  \draw[litedge] (0,-0.5) -- (0.6,-0.5);
  \node[font=\scriptsize, anchor=west] at (0.72,-0.5) {literature (cited)};
  \draw[openedge] (0,-1.0) -- (0.6,-1.0);
  \node[font=\scriptsize, anchor=west] at (0.72,-1.0) {open};
\end{scope}
\end{tikzpicture}
\caption[Compression and learnability.]{The top row is settled and kernel-verified
  (\leanname{fundamental_vc_compression}): PAC learnability, finite VC dimension, and
  the existence of a finite compression scheme with side information are all equivalent.
  This three-way equivalence is the qualitative core; it certifies that a short summary
  exists but says nothing about its length. The vertical descent tracks how short. A
  compression of size $2^{O(d)}$ is known, a result that was open until Moran and
  Yehudayoff established it in 2016; whether the size can reach $\mathrm{poly}(d)$ or
  the conjectured $O(d)$ is unresolved. The right branch asks whether the side information
  can be dropped at the same size, also open. Unlabeled compression of bounded size exists
  for every VC class; whether it reaches $O(d)$ is the Kuzmin--Warmuth question, open in
  general and settled only for maximum classes.}
\label{fig:compression-landscape}
\end{figure}

The following table summarizes the best known compression bounds:

\begin{table}[ht]
\centering\small
\caption{Known compression bounds by class family.}
\label{tab:compression-bounds}
\begin{tabular}{lcc}
\toprule
\textbf{Class family} & \textbf{VC dim.} & \textbf{Compression size} \\
\midrule
Maximum classes & $d$ & $d$ \\
Intersection-closed & $d$ & $d$ \\
Halfspaces in $\R^d$ & $d+1$ & $d+1$ \\
VC dim.\ 1 & 1 & 1 \\
Arbitrary VC class & $d$ & $2^{O(d)}$ \\
Unlabeled (arbitrary) & $d$ & No $f(d)$ bound \\
\bottomrule
\end{tabular}
\end{table}

\begin{openprob}[Polynomial compression]
\label{op:poly-compression}
Does every concept class of VC dimension $d$ admit a compression scheme of
size $\mathrm{poly}(d)$?  Even this weaker form of the compression conjecture
remains open.
\end{openprob}

Tight sample complexity, compression, and Occam's razor all express a single
structural fact: \emph{learnability is equivalent to the existence of a short
summary of the evidence}.  The VC dimension certifies that such a summary
exists.  The compression conjecture asks how short it can be.  After forty
years, that question stands unanswered.


% ================================================================
% OPEN PROBLEMS
% ================================================================

\section{Open Problems}
\label{sec:compression-open}

The $O(d)$ conjecture (\Cref{op:compression}) is the gravitational center, but
the frontier it anchors is wider.  The following ten problems branch off it,
some sharply posed with the obstruction articulated, others still missing the
right invariant or vocabulary.

\begin{openprob}[Compression without side information]
\label{op:side-info}
The existence theorem produces a compression scheme \emph{with} side
information: the reconstruction is handed $O(d)$ extra bits beyond the subsample.
Does every class of VC dimension $d$ admit a scheme of size $2^{O(d)}$ with
\emph{no} side bits at all (\leanname{Open_NoInfoCompressionStrengthening})?  A
positive answer is the first step toward
\Cref{op:compression}~\cite{MoranYehudayoff2016}.
\end{openprob}

\begin{openprob}[Orientations of non-maximum graphs]
\label{op:oig-orientation}
\Cref{thm:one-inclusion-orientation} orients the one-inclusion graph of a
\emph{maximum} class with in-degree $\leq d$.  Does every (non-maximum) class of
VC dimension $d$ admit an orientation of in-degree $\mathrm{poly}(d)$?  An
affirmative answer would yield polynomial compression.  No orientation theorem
is known off the maximum case: the combinatorial object is in hand, the extremal
lever that tames it is not~\cite{KuzminWarmuth2007}.
\end{openprob}

\begin{openprob}[Agnostic compression]
\label{op:agnostic-compression}
The schemes of this chapter are realizable: they reconstruct a hypothesis
\emph{consistent} with the sample.  Does every VC-$d$ class admit an
\emph{agnostic} compression scheme of size $O(d)$, one whose reconstruction has
error at most $\mathrm{opt} + \varepsilon$ when no consistent hypothesis
exists?  Sharpened by David, Moran, and
Yehudayoff~\cite{DavidMoranYehudayoff2016}.
\end{openprob}

\begin{openprob}[Compression for attention classes]
\label{op:attention-compression}
\Cref{prop:attention-compression} certifies that the softmax-attention class
$\mathcal{A}_{d,n_K}$ is compressible at size $2^{O(d)}$
(\leanname{softAttention_wellBehaved}), and the support-key heuristic predicts
$O(n_K)$.  Does $\mathcal{A}_{d,n_K}$ admit a compression scheme of size
$O(n_K)$, or even $O(\VC)$?  This is \Cref{op:compression} instantiated for a
contemporary architecture, with no compression invariant specific to attention
yet known~\cite{MoranYehudayoff2016}.
\end{openprob}

\begin{openprob}[Compression and recursive teaching]
\label{op:rtd}
The recursive teaching dimension $\mathrm{RTD}(\mathcal{H})$ is sandwiched with
the VC dimension but the two can differ.  Is the optimal compression size
$\Theta(\mathrm{RTD})$ rather than $\Theta(\VC)$?  This would re-route the
conjecture through teaching rather than shattering~\cite{SimonZilles2015}.
\end{openprob}

\begin{openprob}[Stable and replicable compression]
\label{op:stable-compression}
Does every VC class admit a compression scheme that is \emph{stable} under
resampling, such that changing one training example changes $O(1)$ of the selected
indices?  Stability would link compression to differential privacy and
replicability, where an equivalence with online learnability is already
known~\cite{BunLivniMoran2020}.  The stability invariant for compression has no
accepted definition yet.
\end{openprob}

\begin{openprob}[List compression and the DS dimension]
\label{op:list-compression}
In multiclass learning, finite Daniely--Shalev-Shwartz (DS) dimension
characterizes learnability.  Is finite DS dimension equivalent to a bounded
\emph{list} compression scheme, where reconstruction outputs a short list of
hypotheses containing a consistent one?  The multiclass compression vocabulary
is still being built~\cite{BrukhimCarmonDinurMoranYehudayoff2022}.
\end{openprob}

\begin{openprob}[Fat-shattering compression]
\label{op:fat-compression}
For $[0,1]$-valued function classes, learnability is governed by the
fat-shattering dimension $\mathrm{fat}_\gamma$.  Does finite $\mathrm{fat}_\gamma$
imply a real-valued compression scheme of size $\mathrm{poly}(\mathrm{fat}_\gamma)$?
The Boolean machinery does not transfer verbatim, and the scale parameter
$\gamma$ has no settled role in the compressed object (Alon, Ben-David,
Cesa-Bianchi, and Haussler).
\end{openprob}

\begin{openprob}[Super-linear lower bounds]
\label{op:compression-lower}
Every known VC-$d$ class compresses to $O(d)$ points; the obstruction to the
conjecture is constructing schemes, not ruling them out.  Is there a VC-$d$
class provably \emph{requiring} compression size $\omega(d)$?  No super-linear
lower bound is known, and a resolution either way would be decisive.
\end{openprob}

\begin{openprob}[Boosting weak compression]
\label{op:weak-compression}
Call a scheme \emph{weak} if its reconstruction is correct on a
$\tfrac{1}{2}+\gamma$ fraction of relabelings of the sample.  Does weak
compression boost to full compression with only $O(\log(1/\gamma))$ overhead in
size, mirroring the boosting of weak PAC learners?  A boosting operator on
compression schemes has not been formulated~\cite{MoranYehudayoff2016}.
\end{openprob}


% ================================================================
% EXERCISES
% ================================================================

\subsection*{Exercises}

\begin{enumerate}[leftmargin=*,itemsep=6pt]

\item \textbf{Compression lower bound.}
  Let $\mathcal{H}$ be a concept class with $\VC(\mathcal{H}) = d$, and
  let $(\kappa, \rho)$ be a sample compression scheme of size $k$ for
  $\mathcal{H}$.  Prove that $k \geq d$.

  \emph{Hint:} Let $C = \{x_1, \ldots, x_d\}$ be a shattered set.
  For each of the $2^d$ labelings $b \in \{0,1\}^d$, the compression
  function selects a subsequence $\kappa(S_b) \subseteq S_b$ of size
  $\leq k$.  Show that if $k < d$, a pigeonhole argument forces two
  distinct labelings $b \neq b'$ to compress to the same subsequence
  $\kappa(S_b) = \kappa(S_{b'})$ with the same labels, contradicting
  the requirement that reconstruction produces hypotheses consistent
  with both $S_b$ and $S_{b'}$.

\item \textbf{Compression for unions of intervals.}
  Let $\mathcal{H}_k$ be the class of unions of at most $k$ disjoint
  intervals on $[0,1]$: $h \in \mathcal{H}_k$ iff $h^{-1}(1)$ is a
  union of at most $k$ intervals.
  \begin{enumerate}[label=(\alph*)]
  \item Prove that $\VC(\mathcal{H}_k) = 2k$.
  \item Construct an explicit compression scheme of size $2k$ for
    $\mathcal{H}_k$.  (\emph{Hint:} Generalize
    \Cref{ex:oig-intervals}; the compressed subsample consists of
    the $2k$ boundary points of the positive regions.)
  \item Use the result of Exercise~1 to conclude that no compression
    scheme of size $< 2k$ exists.  This verifies the compression
    conjecture for $\mathcal{H}_k$ with equality.
  \end{enumerate}

\item \textbf{Compression and teaching dimension.}
  The \emph{teaching dimension} $\mathrm{TD}(\mathcal{H})$ is the
  minimum $k$ such that every concept $h \in \mathcal{H}$ has a
  \emph{teaching set} of size $k$: a set $T \subseteq X$ with
  $|T| = k$ such that $h$ is the unique concept in $\mathcal{H}$
  consistent with $h|_T$.
  \begin{enumerate}[label=(\alph*)]
  \item Show that if $\mathcal{H}$ has teaching dimension $k$, then
    $\mathcal{H}$ admits a compression scheme of size $k$.  (The
    teaching set is the compressed subsample.)
  \item Show that a compression scheme of size $k$ does \emph{not}
    imply teaching dimension $\leq k$: construct a class $\mathcal{H}$
    with a compression scheme of size $1$ but teaching dimension
    $> 1$.  (\emph{Hint:} The compression function may depend on the
    full sample, not just on the target concept; the teaching set must
    work for a single concept without knowing the other data.)
  \item Prove that $\mathrm{TD}(\mathcal{H}) \geq \VC(\mathcal{H})$
    for all classes $\mathcal{H}$.  Combined with the compression
    conjecture, this would imply
    $\VC \leq \text{compression size} \leq \mathrm{TD}$: compression
    sits between VC dimension and teaching dimension.  But whether
    the compression conjecture implies
    $\text{compression size} \leq O(\mathrm{TD})$ or
    $\text{compression size} \leq O(\VC)$ remains open; these are
    distinct (and incomparable) questions.
  \end{enumerate}

\end{enumerate}

\subsection*{Counterfactual exercises}

\Cref{fig:compression-landscape} encodes an argument: each node carries a claim
at a specific evidence grade (kernel-verified, literature-cited, or open).  The exercises below stress-test that structure: take a known result,
vary one of its hypotheses or conclusions, and track exactly where the shifted
claim lands.  Each variation reveals why a rung on the size ladder is a real
boundary, why the side-information qualifier matters, and why the distance from
$2^{O(d)}$ to $O(d)$ has resisted forty years of effort.

\begin{enumerate}[leftmargin=*,itemsep=3pt]
\item \textbf{The qualitative core.} Show that the top row is a single three-way equivalence: PAC learnability, finite VC dimension, and the existence of \emph{some} finite compression scheme with side information (\cref{thm:fundamental}, \leanname{fundamental_vc_compression}). Explain why this asserts \emph{that} a short summary exists but says nothing about its length.

\item \textbf{The easy half.} Prove that any finite compression scheme forces finite VC dimension, with the explicit bound $\VC(\mathcal{H}) \le 2(k+1)^2-1$ (\cref{thm:compression-pac}, \leanname{compression_imp_vcdim_finite}). Contrast this settled direction with the converse, finite VC implying \emph{small} compression, which splinters into the literature bound and the open conjectures.

\item \textbf{Halfspaces hit the target.} For halfspaces in $\R^d$, the maximum-margin hyperplane is determined by at most $d+1$ support vectors (\cref{ex:svm-compression}), so the family compresses to size $d+1$ at $\VC=d+1$, attaining the $O(d)$ target. Argue that this support-vector mechanism has no known analogue for an arbitrary VC class.

\item \textbf{Maximum classes collapse the ladder.} For a maximum class of VC dimension $d$ (growth function meeting Sauer--Shelah with equality), show that compression to exactly $d$ points exists, so the rungs $2^{O(d)}$, $\mathrm{poly}(d)$, and $O(d)$ coincide. The vertical separation between rungs is a statement about \emph{arbitrary} classes; explain why it vanishes here.

\item \textbf{The side-information axis.} The existence theorem produces a scheme \emph{with} side information (\leanname{fundamental_vc_compression}); whether the side bits can be removed at the same size is open (\cref{op:side-info}, \leanname{Open_NoInfoCompressionStrengthening}). Show that the two endpoints differ by a structural qualifier alone, the size held fixed.

\item \textbf{The exponential gap.} The best known bound is $2^{O(d)}$ while the conjecture demands $O(d)$ (\cref{rmk:compression-gap}): at $d=20$ this is a gap from $\sim 10^6$ to $\sim 20$. Explain why the $2^{O(d)}$ bound (\cref{thm:moran-yehudayoff}, \cite{MoranYehudayoff2016}), unknown before 2016, is a literature result rather than a refinement of the conjecture.

\item \textbf{The unlabeled variant.} Bounded-size unlabeled compression exists for every VC class (\cref{def:unlabeled-compression}, \cref{thm:unlabeled-bounded}); whether it reaches the labeled $O(d)$ bound is the Kuzmin--Warmuth question, open in general and settled for maximum classes. Explain what dropping the stored labels costs, and why it leaves the size question open rather than closing it.

\item \textbf{Compression as a dimension functor.} Read ``the smallest summary that reconstructs the data'' as a universal property: the optimal compression size as the minimal generating section of an evidence-functor on finite samples whose rank is the VC dimension. Discuss whether such a single dimension functor exists, unifying the size ladder, the side-information qualifier, and the alternative dimensions (recursive teaching, DS), or whether the notions are provably non-unifiable (\cref{op:dimension-functor}, \cref{op:rtd}).
\end{enumerate}
